%表紙
\newpage
\chapter{講義}
\begin{enumerate}[label=第\arabic*条　]
  \item この章は、憲章の趣旨に基づき、講義の実施に関し必要な事項を定めるものとする。
  \item 学園に、音ゲー基礎学部及び音ゲー実践学部を置く。
  \begin{enumerate}[label=\arabic*　,start=2,leftmargin=*]
    \item 各学部に置く講義区分は、次のとおりとする。
  \end{enumerate}
    \begin{enumerate}[label=(\arabic*)　]
      \item 音ゲー基礎学部　創作系及び文理系
      \item 音ゲー実践学部　アーケード系、スタンドアロン系及びモバイル系
    \end{enumerate}
\item 音ゲー基礎学部は、音楽ゲームをめぐる知識、文化及び創作を扱う。
    \begin{enumerate}[label=(\arabic*)　]
      \item 創作系　芸術、設計、開発その他の制作に関する分野
      \item 文理系　人文科学、社会科学及び自然科学に関する分野
    \end{enumerate}
\item 音ゲー実践学部は、音楽ゲームをプレイすることに関する技術及び知見を扱う。
    \begin{enumerate}[label=(\arabic*)　]
      \item アーケード系　店舗等に設置される筐体を用いる音楽ゲームに関する分野
      \item スタンドアロン系　家庭用ゲーム機、計算機その他の機器で動作する音楽ゲームに関する分野
      \item モバイル系　携帯端末で動作する音楽ゲームに関する分野
    \end{enumerate}
\item 前二条に定める学部及び講義区分は、講義を探すための目安であって、学生の活動の範囲を限定するものではない。
  \begin{enumerate}[label=\arabic*　,start=2,leftmargin=*]
    \item 複数の講義区分にまたがる講義については、主たるものを選べば足りる。
    \item いずれの講義区分にも当てはめにくい講義については、最も近いものを選ぶ。この場合において、教務主事は、講義区分の選択について相談に応じる。
  \end{enumerate}
  \item 講師は、所定の申請を行い、その認可を受けて、講義を開講することができる。
    \begin{enumerate}[label=\arabic*　,start=2,leftmargin=*]
      \item 前項の認可は、教務主事の全員一致による。
      \item 第１項の申請の具体的な手続については、別に定める。
    \end{enumerate}
  \item 学生は、別に定めるところにより、講義を受講することができる。
    \begin{enumerate}[label=\arabic*　,start=2,leftmargin=*]
      \item 前項の規定にかかわらず、講師は、担当する講義について、受講に関する条件を定めることができる。
    \end{enumerate}
  \item 講義資料の著作権その他の知的財産権は、当該資料を作成した講師に帰属する。
    \begin{enumerate}[label=\arabic*　,start=2,leftmargin=*]
      \item 前項の規定にかかわらず、学園は、第６章第２条第１項第１号の同意に基づき、学内活動の目的に必要な範囲において、講義資料を無償で使用し、及び配信することができる。
    \end{enumerate}
  \item 講師は、あらかじめ予定した講義の回数を正規活動期間内に実施することが難しい場合、補講期間内に、その不足する回数を実施することができる。
  \item 教務主事は、講義について次の各号のいずれかに該当する事実があると認めるときに限り、当該講師に対し、その中止又は変更を求めることができる。
      \begin{enumerate}[label=(\arabic*)　]
        \item 明らかに虚偽の情報の伝達その他の講義の真正性を著しく損なう事実があるとき
        \item 受講者その他の学生に対するハラスメントその他の著しく不当な言動があるとき
        \item 前二号に掲げるもののほか、学生の権利利益を著しく損なうおそれがあると認められる事実があるとき
      \end{enumerate}
    \begin{enumerate}[label=\arabic*　,start=2,leftmargin=*]
      \item 前項の求めは、憲章の定める学問の自由を不当に制限することのないよう、必要な最小限度にとどめなければならない。
    \end{enumerate}
  \item この章の実施に関し必要な事項は、教務主事が別に定める。
\end{enumerate}
